% SPDX-License-Identifier: GPL-3.0-or-later
% Copyright (C) 2026 Alessandro Greco
%% Testo alternativo per l'accessibilita' (EPUB).
\altfig{Topologia di test del libro configurata con Chutney sulla sola interfaccia 127.0.0.1: tre autorità di directory test000a, 001a e 002a, anch'esse relay, con ORPort 5100–5102 e DirPort 7100–7102; tre relay intermedi 003r–005r e due exit 006e–007e, un client 008c con SOCKS 9008, un onion service 009h con porta virtuale 5858 inoltrata a 127.0.0.1:4747. Le autorità votano ogni 20 secondi nel laboratorio; un esempio di circuito per traffico HTTP va dal client attraverso 003r e 005r verso l'exit 006e e il server di test 127.0.0.1:8080, consentito solo dalla policy di laboratorio. Il servizio onion è distinto e non usa un exit relay; la figura mostra soltanto alcune porte, la tabella del capitolo le elenca tutte. La rete in loopback non equivale a un isolamento dalla rete host.}
% Rete privata di laboratorio: schema semplificato
\begin{tikzpicture}[font=\sffamily\scriptsize,
  aut/.style={relay,fill=torviola!15,draw=torviola},
  ex/.style={relay,fill=torverde!30}]
  \foreach \i/\x in {0/0,1/2,2/4} \node[aut] (a\i) at (\x,3.2) {00\i a};
  \node[above=1pt of a1,font=\sffamily\scriptsize,align=center] {autorità di directory\\ORPort 5100--5102, DirPort 7100--7102};
  \node[font=\sffamily\scriptsize,text=torgrigio,align=center] at (6.6,3.2) {voto ogni 20 s\\consenso firmato};
  \foreach \i/\x in {3/0,4/2,5/4} \node[relay] (r\i) at (\x,1.2) {00\i r};
  \node[ex] (e6) at (6.6,1.2) {006e};
  \node[ex] (e7) at (6.6,-0.6) {007e};
  \node[client] (c) at (-2.4,1.2) {008c client};
  \node[font=\ttfamily\scriptsize,below=1pt of c] {SOCKS 9008};
  \node[client,fill=torverde!10,draw=torverde!60!black] (h) at (2,-1.2) {009h servizio onion};
  \node[font=\ttfamily\scriptsize,below=1pt of h] {5858 $\to$ 127.0.0.1:4747};
  \node[server] (w) at (9.4,1.2) {\texttt{127.0.0.1:8080}};
  \begin{scope}[on background layer]
    \foreach \n in {r3,r4,r5,e6,e7} \draw[dotted,torgrigio] (\n) -- (a1);
  \end{scope}
  \draw[cellaf] (c) -- (r3);
  \draw[cellaf] (r3) to[bend right=40] (r5);
  \draw[cellaf] (r5) -- (e6);
  \draw[cellaf] (e6) -- (w);
  \node[font=\sffamily\scriptsize,text=torviola] at (2,-0.05) {esempio di circuito HTTP};
  \node[font=\sffamily\tiny,text=torgrigio] at (4,-2.4) {onion: nessun exit};
\end{tikzpicture}
